% BÀI TOÁN HÌNH HỌC - BAI_HINH_22 (TRÍCH ĐỀ 22)
\begin{ex}
Cho điểm A nằm ngoài đường tròn tâm O bán kính R. Từ A kẻ đường thẳng (d) không đi qua tâm O, cắt (O) tại B và C (B nằm giữa A và C). Các tiếp tuyến với đường tròn (O) tại B và C cắt nhau tại D. Từ D kẻ DH vuông góc với AO (H nằm trên AO), DH cắt cung nhỏ BC tại M. Gọi I là giao điểm của DO và BC.
	\begin{enumerate}
		\item[a)] Chứng minh bốn điểm O, H, D, C thuộc một đường tròn.
		\item[b)] Chứng minh:
		\begin{itemize}
			\item[+)] $OH \cdot OA = OI \cdot OD$.
			\item[+)] AM là tiếp tuyến của đường tròn (O).
		\end{itemize}
	\end{enumerate}
\loigiai{
{
	\begin{center}
	\begin{tikzpicture}[line join = round, line cap=round,>=stealth,font=\footnotesize,scale=1]
		% Def
		\def\r{3}
		\def\a{5}
		\pgfmathsetmacro{\alphadeg}{acos(\r/\a)} 
		\def\kparam{5/8}
		
		% Coordinates
		\coordinate (O) at (0,0);
		\coordinate (D) at (\a,0);
		\coordinate (C) at (\alphadeg:\r);
		\coordinate (B) at (-\alphadeg:\r); 
		
		\coordinate (I) at ($(B)!0.5!(C)$);
		
		\coordinate (A) at ($(B) - \kparam*($(C)-(B)$)$);
		
		\coordinate (H) at ($(O)!(D)!(A)$);
		
		\path [name path = c1] (O) circle (\r);
		\path [name path = oa] (O)--(A);
		\path [name intersections={of=c1 and oa, by={K}}];
		
		\path [name path = dh] (D)--(H);
		\path [name intersections={of=c1 and dh, by={M}}];
		
		% Draw 
		\draw (O) circle (\r); 
		
		\draw (O) -- (C) -- (D) -- (B) -- cycle;
		\draw (D) -- (O) -- (A) -- (C);
		\draw (D) -- (H);
		\draw[dashed] (O) -- (M) -- (A);
		\draw[dashed] (M) -- (K);
		
		% Fill points (theo hướng dẫn của bạn)
		\fill (O) circle (1.5pt);
		\fill (D) circle (1.5pt);
		\fill (C) circle (1.5pt);
		\fill (B) circle (1.5pt);
		\fill (I) circle (1.5pt);
		\fill (A) circle (1.5pt);
		\fill (H) circle (1.5pt);
		\fill (K) circle (1.5pt);
		\fill (M) circle (1.5pt);
		
		% Node points (theo hướng dẫn của bạn)
		\node[below left] at (O) {$O$};
		\node[below right] at (D) {$D$};
		\node[above right] at (C) {$C$};
		\node[below right] at (B) {$B$};
		\node[below] at (I) {$I$};
		\node[below] at (A) {$A$};
		\node[left] at (H) {$H$};
		\node[below left] at (K) {$K$};
		\node[above right] at (M) {$M$};
		
	\end{tikzpicture}	
		
	\end{center}
}
\end{ex}
